\begin{trapbox}
	\begin{description}[style=nextline, leftmargin=2.8em, labelsep=0.5em, itemsep=0.35em]
		\item[\textbf{\textcolor{CTrap}{易错点一（随意相消）}}] 在数列求和中，误认为裂项后任意$k$都能相邻项相消。
		\item[\textbf{\textcolor{CTrap}{正确理解（k=1可消）：}}] 只有$k=1$时，$\sqrt{n+1}-\sqrt{n}$与下一项$n+1$的表达式才直接相消；一般$k$不能直接相邻相消。
		\item[\textbf{\textcolor{CTrap}{错因分析（间隔误解）：}}] 对裂项结构未分析透彻，忽略$n$和$n+k$的间隔。
	\end{description}

	\medskip
	\begin{description}[style=nextline, leftmargin=2.8em, labelsep=0.5em, itemsep=0.35em]
		\item[\textbf{\textcolor{CTrap}{易错点二（共轭方向错）}}] 分子分母同乘时，错误地乘以$\sqrt{n+k}+\sqrt{n}$，而非其共轭。
		\item[\textbf{\textcolor{CTrap}{正确理解（乘共轭差）：}}] 必须乘以$\sqrt{n+k}-\sqrt{n}$才能使用平方差公式；乘以和会导致分母平方，无法化简。
		\item[\textbf{\textcolor{CTrap}{错因分析（平方差误）：}}] 混淆了平方差公式需“和乘差”的结构。
	\end{description}

	\medskip
	\begin{description}[style=nextline, leftmargin=2.8em, labelsep=0.5em, itemsep=0.35em]
		\item[\textbf{\textcolor{CTrap}{易错点三（漏系数k）}}] 直接写$\frac{1}{\sqrt{n+k}+\sqrt{n}}=\sqrt{n+k}-\sqrt{n}$，忘记除以$k$。
		\item[\textbf{\textcolor{CTrap}{正确理解（分母含k）：}}] 公式右端为$\frac{\sqrt{n+k}-\sqrt{n}}{k}$，分母有$k$，不可省略。
		\item[\textbf{\textcolor{CTrap}{错因分析（常数漏除）：}}] 平方差结果为$k$没有移分母，错误认为结果为1。
	\end{description}
\end{trapbox}
